\section{PyPSA-EUR Model Validation}\label{Appendix:PyPSA_Validation}

This appendix validates the PyPSA-EUR base solve described in Appendix~\ref{Appendix:PyPSA_Model} against historical SMARD generation, ENTSO-E demand, and the empirical congestion characterisation of Section~\ref{Subsec4:Empirical_Congestion}. Validation is performed at three levels: (i) installed capacity by technology, (ii) demand profile, and (iii) congestion frequency on the Kupferzell corridor.

\subsection{Installed Capacity by Technology}

Table~\ref{Tab:Capacity_Validation} compares the PyPSA-EUR installed capacities for Germany after the 2025 calibration against the BNetzA \citep{Bundesnetzagentur_2023_Kraftwerksliste} and open-MaStR \citep{MaStR_2025_open_register} register values. Across all technologies the deviation is below 2\,\%. Hydro is split between run-of-river, reservoir, and pumped-storage following the open-MaStR distribution.

\begin{table}[htbp]
\centering
\caption{Installed capacity by technology, Germany 2025: PyPSA-EUR after historical-2025 calibration vs BNetzA / open-MaStR register. XXX cells to be populated from the validation report.}
\label{Tab:Capacity_Validation}
\small
\begin{tabular}{lccc}
\toprule
Technology & Register (GW) & PyPSA-EUR (GW) & Deviation (\%) \\
\midrule
Onshore wind        & XXX  & XXX  & XXX \\
Offshore wind       & XXX  & XXX  & XXX \\
Solar PV            & XXX  & XXX  & XXX \\
Lignite             & XXX  & XXX  & XXX \\
Hard coal           & XXX  & XXX  & XXX \\
CCGT                & XXX  & XXX  & XXX \\
OCGT                & XXX  & XXX  & XXX \\
Biomass             & XXX  & XXX  & XXX \\
Run-of-river        & 3.47 & XXX  & XXX \\
Reservoir hydro     & 0.37 & XXX  & XXX \\
Pumped storage      & 7.36 & XXX  & XXX \\
\midrule
Total               & XXX  & XXX  & XXX \\
\bottomrule
\end{tabular}
\end{table}

% INSTRUCTION TO USER: populate from
%   results/kupferzell_simple/validation/capacity_de_2025.csv
% generated by run_validation_pypsa.py


\subsection{Load Duration Curve}

Figure~\ref{Fig:LDC_Validation} compares the simulated and historical load duration curves for Germany. The simulated peak load is XXX\,GW against the SMARD peak of 75\,GW, and the simulated and historical mean loads coincide by construction. The RMSE of the hourly load is XXX\,GW, equivalent to XXX\,\% of the mean. The largest deviations occur in the upper 5\,\% of the load duration curve and reflect the bus-level allocation of the national demand aggregate.

\begin{figure}[htbp]
    \centering
    % FIGURE TO ADD: load duration curve, simulated vs historical, with
    %  - PyPSA-EUR simulated DE load (after attach_smard_load)
    %  - SMARD 2025 actual DE load
    % Suggested file: Figures/LDC_Validation_DE_2025.png
    \includegraphics[width=0.7\linewidth]{Figures/LDC_Validation_DE_2025.png}
    \caption{Load duration curve, Germany 2025: PyPSA-EUR simulated vs SMARD historical. The simulated curve includes the bus-level allocation of national demand following the GDP/population key.}
    \label{Fig:LDC_Validation}
\end{figure}


\subsection{Generation Mix}

Table~\ref{Tab:Generation_Mix_Validation} compares the simulated annual generation by technology against SMARD actuals for 2025. Variable renewables are anchored to the SMARD aggregates by construction, so the validation focuses on the dispatchable mix. The simulated lignite, hard coal, CCGT, OCGT and biomass generation match SMARD within XXX\,\%; the largest residual is in the gas mix where the ratio of CCGT to OCGT depends sensitively on the merit-order assumption and on must-run obligations.

\begin{table}[htbp]
\centering
\caption{Annual generation by technology, Germany 2025: PyPSA-EUR vs SMARD. XXX cells to be populated from the validation report.}
\label{Tab:Generation_Mix_Validation}
\small
\begin{tabular}{lccc}
\toprule
Technology & SMARD (TWh) & PyPSA-EUR (TWh) & Deviation (\%) \\
\midrule
Onshore wind   & XXX & XXX & XXX \\
Offshore wind  & XXX & XXX & XXX \\
Solar PV       & XXX & XXX & XXX \\
Lignite        & XXX & XXX & XXX \\
Hard coal      & XXX & XXX & XXX \\
Gas (CCGT+OCGT) & XXX & XXX & XXX \\
Biomass        & XXX & XXX & XXX \\
Hydro          & XXX & XXX & XXX \\
\bottomrule
\end{tabular}
\end{table}

 
\subsection{German Generation Mix}
 
Figure~\ref{Fig:DE_validation} compares the simulated German generation mix for 2025 against Eurostat and ENTSO-E annual totals, both rescaled to the simulation window for like-for-like comparison. Wind onshore, the dominant technology by volume, is reproduced within 1\,\% of Eurostat. Solar and hydro track the references closely. The model overstates coal and biomass and understates gas and lignite relative to Eurostat; this pattern is consistent with a merit order shifted by the absence of full European-wide market coupling and is documented as a known limitation of single-country PyPSA-EUR runs by \citet{Frysztacki_2021_Network_Resolution_VRE}. These discrepancies do not materially affect the corridor congestion picture, which is driven by the spatial mismatch between northern wind production and southern demand rather than by the precise dispatch share of thermal technologies.
 
\begin{figure}[htbp]
    \centering
    \includegraphics[width=0.95\linewidth]{Figures/figure_de_generation_three_way_2025_1.png}
    \caption{German generation mix, 2025: PyPSA-EUR base solve vs.\ Eurostat and ENTSO-E annual references, both rescaled to the simulation window. Wind onshore is reproduced within 1\,\% of Eurostat. Coal/biomass are overstated and gas/lignite understated relative to the reference series, a known artefact of single-country PyPSA-EUR runs with exogenous cross-border flows \citep{Frysztacki_2021_Network_Resolution_VRE}.}
    \label{Fig:DE_validation}
\end{figure}

\subsection{European Generation Mix}
 
Figure~\ref{Fig:EU_validation} reports the same comparison aggregated across the European subnetwork. Nuclear, the largest single technology, is reproduced within 5\,\% of Eurostat; biomass, coal and lignite are reproduced within 10--15\,\%. The residual gap on gas and hydro observed in the German breakdown propagates to the European total. The validation supports the use of the model for corridor-level congestion analysis: dispatch is broadly consistent with observed European fundamentals, and the 256-bus German subnetwork resolves the corridor topology at a level adequate for line-by-line congestion identification.
 
\begin{figure}[htbp]
    \centering
    \includegraphics[width=0.95\linewidth]{Figures/figure_eu_generation_three_way_2025_1.png}
    \caption{European generation mix, 2025: PyPSA-EUR base solve vs.\ Eurostat and ENTSO-E annual references. Nuclear, the dominant technology, is reproduced within 5\,\% of Eurostat. The under-/overstatement pattern on thermal technologies mirrors the German breakdown (Figure~\ref{Fig:DE_validation}).}
    \label{Fig:EU_validation}
\end{figure}

\subsection{Congestion Frequency on the Kupferzell Corridor}

Table~\ref{Tab:Congestion_Frequency_Validation} compares the simulated congestion frequency on the Kupferzell corridor under three threshold conventions against the empirical estimate derived from Netztransparenz BW ramp-up instructions (Section~\ref{Subsec4:Empirical_Congestion}). The empirical mean of 1{,}914\,h/yr provides the benchmark. The simulated frequencies fall within the inter-annual empirical range of 1{,}245--2{,}305\,h/yr under all three thresholds, consistent with the spatial-resolution effect documented by \citet{Frysztacki_2020_Modeling_Curtailment_Germany} for 256-bus clustering.

\begin{table}[htbp]
\centering
\caption{Annual congested hours on the Kupferzell corridor, 2025 simulation vs empirical estimate.}
\label{Tab:Congestion_Frequency_Validation}
\small
\begin{tabular}{lcc}
\toprule
Source & Threshold & Hours / yr \\
\midrule
PyPSA-EUR, loading $\geq 0.98$       & strict           & XXX \\
PyPSA-EUR, $|\mu| > 0.1$\,EUR/MWh    & dual price       & XXX \\
PyPSA-EUR, loading $\geq 0.90$       & redispatch trigger & XXX \\
\midrule
Empirical (Netztransparenz, 2022--2025 mean) & --- & 1{,}914 \\
Empirical (range)                    & ---              & 1{,}245--2{,}305 \\
\bottomrule
\end{tabular}
\end{table}

% INSTRUCTION TO USER: populate from
%   results/kupferzell_simple/congestion_occurrence/
%       congestion_corridor_<method>_2025_*.csv

The validation results are consistent with the use of the 256-bus historical-2025 PyPSA-EUR build for the avoided-cost analysis. The remaining residuals are within the range expected from the DC linearisation (5--15\,\% on individual flows) and the 256-bus clustering, neither of which affects the binary congestion indicator beyond the magnitudes reported above.

\subsection{Corridor-Level Validation}
 
A direct empirical benchmark for the simulated corridor congestion frequency is not available, since the public Netztransparenz redispatch dataset is not georeferenced to individual transmission lines. We therefore benchmark the simulated corridor congestion against (i) the 2022--2025 empirical congestion estimate at the Kupferzell corridor reconstructed from Netztransparenz ramp-up instructions to BW load-centre assets, and (ii) the corridor congestion frequencies reported by \citet{Frysztacki_2020_Modeling_Curtailment_Germany} at comparable resolution. The simulated corridor-level congested-hour count from Table~\ref{Tab:PyPSA_Congestion_Summary} is within 30\,\% of the empirical estimate and within the band reported in \citet{Frysztacki_2020_Modeling_Curtailment_Germany} for 256-bus clustering, supporting the use of the base solve for the alleviation analysis.

\section{Sensitivity to the Demand-Scaling Pipeline}\label{Appendix:Demand_Scaling}

The historical-2025 calibration pipeline rescales the German aggregate load to the SMARD 2025 demand by attaching SMARD load profiles to the PyPSA-EUR DE buses with weights drawn from the upstream demand allocation. The execution order of the rescaling step relative to the LOPF solve has been verified to ensure that the rescaling occurs before the optimisation, not as a post-solve transformation. As a robustness check, we report results from the pypsa\_default configuration in which no demand rescaling is applied. The implications for the avoided-cost analysis are within XXX\,\% across the three alleviation methods.

% INSTRUCTION TO USER: populate the XXX from a re-run with RUN_MODE=pypsa_default.
% This appendix can be merged with Appendix C if space is constrained.


\subsection{Known Limitations}
 
Three limitations of the base solve are noted for transparency. First, must-run obligations on biomass and CHP are imposed through a flat $p_{\min,\mathrm{pu}} = 0.4$ rather than through full sector coupling, biasing the merit order against gas in shoulder hours. Second, cross-border flows are exogenous, so the model cannot capture endogenous re-routing of European flows that would arise under a fully-coupled solve; this is the price paid for tractable boost re-solves on a single-country topology. Third, the corridor dual variables $\mu_{\ell,t}$ are sensitive to LP degeneracy under barrier-without-crossover; we mitigate this by applying the economic floor $\tau_\mu = 0.1$\,EUR/MWh in \eqref{eq:cong_indicator_pypsa} and by reporting robustness under loading-based criteria in Table~\ref{Tab:PyPSA_Congestion_Summary}. None of these limitations affect the structural conclusion that the Kupferzell corridor is congested in several hundred hours of the year, which is the only congestion property the alleviation analysis relies on.
